\chapter{Glossary}
\label{app:glossary}

\begin{description}[style=nextline,leftmargin=0pt]
\item[Criticality] The state $\keff=1$ in which neutron production balances loss
and the chain reaction is self-sustaining at constant power.
\item[Delayed neutrons] The small fraction ($\beta$) of fission neutrons emitted
seconds to minutes after fission by the decay of precursor fission products; they
set the controllable timescale of a reactor.
\item[Doppler coefficient] The (prompt, negative) change in reactivity per unit
fuel temperature, arising from broadening of $^{238}$U resonances.
\item[Dollar] A unit of reactivity equal to $\beta$; one dollar marks prompt
criticality.
\item[Decay heat] Thermal power from radioactive decay of fission products,
several per cent of full power just after shutdown, declining slowly over days.
\item[Inhour equation] The relation between a constant reactivity and the reactor
period, derived from point kinetics.
\item[Operating reactivity margin (ORM)] The RBMK's shutdown reserve, expressed as
an equivalent number of inserted control rods; a minimum was mandated and was
grossly violated at Chernobyl.
\item[Point kinetics] The reduction of reactor dynamics to ordinary differential
equations for the spatially-averaged power and precursor concentrations.
\item[Positive scram effect] The RBMK phenomenon in which beginning to insert the
graphite-tipped control rods added positive reactivity in the lower core.
\item[Prompt criticality] The state $\rho\ge\beta$ in which prompt neutrons alone
sustain growth; the response collapses to the prompt timescale.
\item[Reactivity] The fractional departure from criticality, $\rho=(\keff-1)/\keff$.
\item[Void coefficient] The change in reactivity per unit coolant void (steam)
fraction; negative in under-moderated water reactors, strongly positive in the
RBMK.
\item[Xenon-135] A fission product with an enormous neutron-absorption
cross-section whose hours-long dynamics cause post-shutdown poisoning and spatial
oscillations.
\end{description}
